\chapter{离线认知}

\section{宿主提供维护时段，Nous 选择认知工作}

维护时段由 Heptalogos 或其他宿主通过维护租约授权，并明确可用时长、资源、网络与模型服务权限、预算和取消条件。Nous Wave 维护待处理的认知需求，并在授权时段内选择经历片段或日记维护、重新验证、梦境探索或服务投影重建。\cite{nous-research-package}

\begin{ReportWideFigure}
\centering
\begin{tikzpicture}[
  node distance=0.5cm,
  block/.style={draw=NWBlue,rounded corners,align=center,text width=3.0cm,minimum height=1.0cm,inner sep=5pt},
  owner/.style={draw=NWRule,rounded corners,align=center,text width=3.8cm,minimum height=1.35cm,inner sep=5pt},
  arrow/.style={-{Stealth[length=2mm]},thick,color=NWBlue}
]
\node[block] (host) {宿主授予维护租约\\预算与取消条件};
\node[block,right=of host] (snapshot) {固定来源快照\\修订、访问控制与序号水位};
\node[block,right=of snapshot] (explore) {有界回放\\关联扩展与检索};
\node[owner,below=1.2cm of explore] (hypothesis) {假设\\支持证据\\反例证据\\适用范围};
\node[owner,left=of hypothesis] (reground) {重新核验来源\\确认当前版本};
\node[owner,left=of reground] (owner) {目标领域所有者\\接受、拒绝或标记冲突};
\draw[arrow] (host) -- (snapshot);
\draw[arrow] (snapshot) -- (explore);
\draw[arrow] (explore) -- (hypothesis);
\draw[arrow] (hypothesis) -- (reground);
\draw[arrow] (reground) -- (owner);
\end{tikzpicture}
\captionof{figure}{离线认知的维护授权、来源快照、假设生成、重新核验与领域路由。假设支持度依据来源证据和反例记录评估。}
\label{fig:offline-cognition}
\end{ReportWideFigure}

\section{梦境探索是有预算的搜索}

梦境探索（Dream）针对有预算的回放与重组。一次可复核的运行固定 SubjectId、输入对象及修订引用、来源快照、访问控制与生命周期、生成者签名、流程版本、预算和权威序号水位。允许的操作包括回放、跨时间比较、反例搜索、冲突核查、有界拓扑扩展和经历片段重组；输出先成为假设（Hypothesis）。\cite{nous-research-package}

假设至少包含主张及范围、输入修订、支持证据、反例证据、不确定性、目标领域、拟议操作、生成者和运行标识。提交前重新检查来源是否存在且可访问、主张是否受来源区域支持、派生谱系是否去重、反例搜索结果，以及当前对象纪元/版本、访问控制、生命周期和清除状态。

\section{停止条件与可靠性}

收益偏低、循环重复、缺少支持的假设比例过高、预算耗尽、目标领域无法处理、前台任务抢占、访问控制或清除状态变化，以及显式取消，都可触发停止。流程记录可解释的停止原因和运行状态。旧快照产生的假设保留原始来源版本；重新核验当前状态后，再由领域所有者决定处理方式。\cite{nous-research-package}

\section{当前实现状态}

研究资料包定义了梦境探索、假设复核和领域路由的目标机制。当前两个代码库分别提供记忆整理和产品运行时能力；跨系统维护租约、反例搜索、重新核验与多领域提案流程列入后续实现。
